\subsubsection{\centering Existence and Uniqueness of Steady States}
\begin{theorem}\label{thm:existence_uniqueness}
	The steady-state solution of the difference equation
	dynamics~\eqref{eq:error_dynamics} exists and is unique for any $\beta_i$,
	$i = 1,\ldots,n$~\cite{tymoshchuk2024tracking}.
\end{theorem}

\begin{proof}
	A steady state satisfies $e(k+1) = e(k)$; substituting into~\eqref{eq:error_dynamics}
	yields the equilibrium condition
	\begin{equation}
		\beta\,\operatorname{sgn}(e^{*}) = 0.
		\label{eq:equilibrium_condition}
	\end{equation}
	Since $\beta_i > 0$ for every $i$, equation~\eqref{eq:equilibrium_condition} holds
	if and only if $\operatorname{sgn}(e_i^{*}) = 0$ for each component $i$.
	By~\eqref{eq:error_dynamics}, this is satisfied uniquely at $e_i^{*} = 0$, so the
	unique steady state is $e^{*} = 0$, equivalently $x^{*}(k) = r(k)$.

	\medskip
	\noindent\textit{Existence.}\enspace
	The cost function $C(x) = \|e\|$ is convex in $e$, which guarantees that its
	infimum is attained.  The gradient update $\Delta e = -\nabla C(e) = -\operatorname{sgn}(e)$
	used in~\eqref{eq:error_dynamics} drives $e$ toward the minimizer $C^{*} = 0$,
	confirming that the equilibrium $e^{*} = 0$ is
	attainable~\cite{tymoshchuk2024tracking}.
\end{proof}
